\chapter{향후 계획}

\section{검증 과제의 완료 상태}

중간보고서에서 정한 검증 과제(영교차 게이트 방향 반전, 선행기술 검색, Stage A 4모듈 경부하 운전과 공차, 보호 상태기계–SALS 연동, 부품 정격 표)는 모두 끝냈다. 부품·설계 판정 표(16장)에 남은 불통과가 없다는 것은 선정 부품이 해당 계산 조건에서 통과했다는 뜻이며, 회로와 평균모델의 불일치(TB33 오동작 사례)와 확인하지 않은 반복 오동작·강제 끔 과도가 남아 있다(12장). 하드웨어 단계의 확인 과제는 표 \ref{t:plan}에 정리했다.

\Needspace{12\baselineskip}\section{하드웨어 단계에서 확인할 설계 요구와 조건부 부품}

{\footnotesize
\begin{longtable}{L{\dimexpr 0.075\linewidth-2\tabcolsep\relax}L{\dimexpr 0.185\linewidth-2\tabcolsep\relax}L{\dimexpr 0.37\linewidth-2\tabcolsep\relax}L{\dimexpr 0.37\linewidth-2\tabcolsep\relax}}
\caption{하드웨어 단계의 확인 과제}\label{t:plan}\\
\toprule
\tabhead \textbf{순서} & \textbf{과제} & \textbf{필요성} & \textbf{완료 기준} \\
\midrule
\endfirsthead
\toprule
\tabhead \textbf{순서} & \textbf{과제} & \textbf{필요성} & \textbf{완료 기준} \\
\midrule
\endhead
\bottomrule
\endlastfoot
1 & 새 인터페이스 요구의 실제 응답 & Stage C 보호 일시 정지와 이중 고장 400 V 부하 정지(값은 표 \ref{t:req})는 모델의 가정이다 & 조정형 Stage C 제어기의 입력 제한 구현, 차량 400 V 부하(V2L·컴프레서·히터)의 정지 응답 측정 \\
2 & V\textsubscript{A} 클램프 문턱 구현 요구 & 여유가 작고(상한 쪽 0.037 V) 켬·끔 지연 1.0 µs 고정에서 정한 요구이며, 0.5 µs에서는 151.26 °C 반례가 있다(12.6절, E44·E45). 최대 연속 켜짐 시간(0.5 ms)의 실제 동작도 미확인 & 검출 회로 설계와 온도·노화 오차 예산, 오동작·잡음 시험. 실제 지연 범위와 판정 기준을 먼저 기록한 새 시험에서 두 지연을 독립 입력으로 사건·가지저항·문턱 오차·폭·수동소자 공차(독립 공차 32개 끝값과 경계 인근)와 교차해 다시 검사 \\
3 & 조건부·단서 부품 & 클램프 저항 891SP의 열용량·저항값별 허용치, 제동 저항 51.9 J 펄스열의 열누적, C\textsubscript{fly} 후보의 다른 온도 전류 정격, 클램프 스위치의 예비 데이터시트·열 기준점 차 & 제조사 확인 또는 부품 시험 \\
4 & 충전 램프와 사전 기동 & 사전 기동과 600 µs 램프를 실제 충전기 응답으로 만족시킬 수 있는지. 640 V·0–1 kW 기동은 공칭 한도 초과(8.3절, E46) & 충전기 종류별 요청–응답 측정, 보조부하·온도·초기 모듈 수에 따른 기동·감발 행렬 \\
5 & 하드웨어 실증 & 실제 지연·잡음·열·절연·전자기 간섭, 공심 L\textsubscript{r}의 누설 자장, 손실 모델(측정 2:1 하드웨어와의 차) & HIL, 축소전력 시제품 \\
6 & Stage C 일시 정지 동안의 48 V & Stage C가 최대 12.4 ms 멈추는 동안 48 V는 배터리 전압을 따른다. 배터리 등가 계산(E43, 10.4절)만 있다 & 48 V 배터리(내부저항·인덕턴스·온도·SOC)와 Stage C 재개 과도를 넣은 결합 시험, 배터리 요구(10.4절)와 48 V 회생 시 충전 수용·BMS 동작 확인 \\
7 & H2 지연 범위 & +3 µs 복합 공차 스위칭 114/128(5장, E46) & 실제 센서·게이트 지연 범위를 정한 뒤 제어기 재조정과 복합 공차 재검사 \\
8 & Stage C CV 설정 범위 & 52 V·270 V·−650 W 단측 송신 듀티 0.980 > 0.97(10.2절, E46) & 36–52 V 설정 전체에서 두 모드 이력을 포함한 정상·과도 점검 \\
9 & W410 자기소자와 최종 클램프 결합 회로 & W410(L\textsubscript{dab} 11.662 µH) 1차 첨두 약 20.35 A는 평균·정적 환산값이고 코어 포화는 미확인. 최종 클램프 + 정격 부하 + 30 ms 충전기 정지를 함께 넣은 회로는 없다(A33·B33은 무부하, 정격 부하 경로 상실은 이전 1 Ω 클램프의 TB28, 12.6절) & W410 변압기 설계·포화 확인, 최종 클램프·부하·30 ms 정지를 함께 넣은 회로 시험 \\
\end{longtable}
}

\section{후속 과제}


\begin{itemize}
\item \textbf{A25R 재확인(18.4절).} 멈춘 모듈 C\textsubscript{fly} 전압 요구, Stage B 보호·강등 시험(E17, E19–E21), 225 kW 확장 회로, E16 사례 집합의 결합 모델 재계산.
\item \textbf{고전압 팩(224·240s) 충전 중 정지}의 Stage B 결합 동작과 회로 확인(지금은 511/506 V 설정의 V\textsubscript{A} 수준 결과만 있음).
\item \textbf{공차와 변형.} C\textsubscript{fly} 후보의 위상 간 편차 최악 조합과 온도에 따른 용량·손실 재계산, Stage A 영전압 스위칭 보조 운전 변형, 대규모 몬테카를로.
\item \textbf{클램프 오동작과 과도.} 오동작 원인(문턱 오차, V\textsubscript{A} 과도·잡음)의 모의, 반복 오동작, 강제 끔 순간의 di/dt·과전압, 0.5 ms 켜짐 때의 Stage A 출렁임(진단 위상 첨두 155.2 A).
\item \textbf{비교와 설계 최적화.} 견인 모터 상저항 범위와 모터 승압 비교의 재대조(경계 4.0–5.2 mΩ), 연속 제약 지령표로 DAB 전이 여유 5 \% 회수, R2-W410 회로 확인, DAB 고장 강등·켬·끔의 통합 회로 확인, Stage A 모듈 순환(마모 분산).
\item \textbf{운전 순서.} 이중 고장 뒤 재기동(V\textsubscript{A} 사전충전 경로, 48 V 배터리 유지 시간, 400 V 부하 재연결), SALS 메시지의 실제 구현(지연 1–5 ms, 재송신 10 ms, 워치독 50 ms는 할당값), 링크 프리차지.
\end{itemize}
